\section{Вступ}
\label{sec:intro}

Permissioned візантійсько відмовостійкий (BFT) консенсус лежить в основі систем
електронного урядування --- голосування, реєстри, цифрова ідентичність --- і
ставить конкретне питання закупівлі: \emph{скільки валідаторів має мати мережа?}

Кількість валідаторів~$\n$ діє в протилежних напрямках. Пропускна здатність
схиляє до малого~$\n$: кожен валідатор виконує кожну транзакцію, а протоколи
родини PBFT обмінюються $O(\n^{2})$ повідомленнями на блок. Надійність
виявлення відмов схиляє до великого~$\n$. У цій статті розвинено відтворювану
інженерну методику, що поєднує обидва чинники за ресурсно-нормованого
калібрування.

Практична перешкода полягає в тому, що виділені багатохостові стенди часто
недоступні на етапі проєктування. Типові CI-раннери доступні й відтворювані, але
розміщують багато валідаторів на невеликій кількості ядер. Без нормування
підігнані показники змішують вартість консенсусу з конкуренцією за ресурси хоста
та з розкладом конкурентності генератора навантаження. Опубліковані
однофакторні підгонки $\TPS(\n)\propto\n^{\alpha}$ тоді розходяться за величиною
і іноді за знаком. Наші попередні дослідження, екстрапольовані до шістдесяти
валідаторів, відрізняються приблизно на порядок --- це часто трактують як
поведінку, специфічну для протоколу. У досліджених умовах показано, що така
розбіжність може виникати лише з дизайну вимірювання.

\paragraph{Внески.}
Центральним методологічним внеском є ресурсно-нормоване вимірювання (\RNM{}):
фіксована CPU-квота~$\q$ на валідатор і $\TPSn=\TPS/\q$, що калібрує вартість
консенсусу окремо від конкуренції за ресурси хоста; у дослідженому розгортанні
воно стабілізує показник за фіксованої квоти, хоча показник і далі залежить від
рівня квоти (розділ~\ref{sec:exp}). Аналітично, за ненасиченої двофакторної
моделі з шумом нульового середнього
$\ahat_{m}\xrightarrow{\Prob}\gamma\lambda-\beta$ вздовж шляху конкурентності
(Пропозиція~\ref{prop:identifiability}) --- це стандартне зміщення через
пропущену змінну~\cite{wooldridge2010}, яке тут використано для інтерпретації
виміряних однофакторних показників і для обґрунтування факторного дизайну
\RNM{} ($\lambda=0$). Також сформульовано задачу з імовірнісними обмеженнями з
допустимою множиною $[\nmin,\nmax]$ (Пропозиція~\ref{prop:window}) і
відкалібровано робочий процес на публічній Besu/QBFT commodity-CI кампанії
(розділи~\ref{sec:exp}--\ref{sec:case}). Компроміс «пропускна здатність ---
детекція» сформульовано, але не продемонстровано: виміряний детектор близький
до випадкового, детекційне обмеження з ним невиконуване, і вікно визначає лише
пропускна здатність (підрозділ~\ref{sec:detection}). Коефіцієнти специфічні
для розгортання і розраховані на перекалібрування; крос-клієнтна реплікація
поза межами.
